\subsection{2026-09-21 --- Feature iteration round 2, and a data-usage audit}
\label{sec:feature-iter-2}

Fifteen builders, 18 cumulative stages, screened on one fit (2015--2019) scored against
2020/2021/2022, with a full rolling-origin confirmation at the end. Registry now lives in
\texttt{src/modeling/pipeline.py} so stages are data, not notebook cells.

\subsubsection*{Results.}

\begin{table}[htbp]
  \centering
  \caption{Round-2 stages, screening scale. ``Up'' = folds improved out of 3.}
  \label{tab:feature-iter-2}
  \begin{tabular}{llrr l}
    \toprule
    Stage & & $\Delta$ & Up & Verdict \\
    \midrule
    06 & price $-$ water value spread & $+0.077$ & 3/3 & confirmed winner \\
    02 & case-relative price          & $+0.036$ & 3/3 & confirmed winner \\
    07 & terminal water value         & $+0.033$ & 3/3 & confirmed winner \\
    01g& \textbf{seasonal, GLOBAL}    & $+0.011$ & 3/3 & \textbf{new winner} \\
    05 & water-value structure        & $+0.010$ & 3/3 & winner \\
    14 & graph cascade (multi-hop)    & $+0.009$ & 2/3 & marginal \\
    09 & inflow regime                & $+0.009$ & 2/3 & marginal \\
    17 & calendar                     & $+0.001$ & 1/3 & nothing \\
    15 & merit order                  & $+0.000$ & 2/3 & nothing \\
    03 & price windows                & $+0.001$ & 2/3 & nothing \\
    06x& \textbf{spread extensions}   & $-0.000$ & 2/3 & \textbf{nothing} \\
    20 & \textbf{water budget}        & $-0.001$ & 1/3 & \textbf{nothing} \\
    04 & price shape                  & $-0.006$ & 0/3 & harmful \\
    10--12 & physics                  & $-0.012$ & 0/3 & harmful \\
    01 & \textbf{seasonal, expanding} & $-0.015$ & 0/3 & \textbf{harmful} \\
    \bottomrule
  \end{tabular}
\end{table}

\finding{\textbf{Final full-CV score 0.8144, up from 0.8080 --- $+0.006$ for 165 additional
features.} A poor return. The first batch (spread, terminal water value, case-relative price) did
essentially all the work; the second batch of nine builders contributed almost nothing
collectively.}

\finding{\textbf{Causality was the problem with the seasonal feature, not the feature.} The
expanding-window version lost 0.015 on 0/3 folds; the \emph{global} version, differing only in
that it ranks against all years rather than prior years, gained 0.011 on 3/3. Early cases have
almost no prior history to rank against, so the expanding version fed the model noise for the
first several years of training. This came directly from the observation that the competition
withholds only test-period \emph{targets}, so feature-side statistics may span the dataset
(Section~\ref{sec:eda-round-3b}). A clean A/B, and it reverses the round-1 rejection.}

\finding{\textbf{Two of my predictions failed outright, and both failures were predicted by the
opposite argument.} (i) The water-budget family --- 41 columns encoding available energy versus
profitable hours, and a greedy merit-dispatch margin --- was argued to be the most promising idea
available on the grounds that it was both hourly-varying and an encoded decision rule, the two
properties that won in round 1. Result: $-0.001$, 1/3 folds. (ii) Twenty-six additional spread
variants, built because the spread family held 60.8\% of model gain, returned $-0.0003$. The
objection raised at the time --- that if the information is already encoded, re-encoding it adds
nothing, and that breadth into unexplored directions matters more than depth on a known winner ---
was correct on both counts.}

\finding{\textbf{Gain share remains a poor guide.} The physics stage held 13.5\% of model gain in
round 1 and measured $+0.001$ there; under the cleaner screening comparison it measures
$\mathbf{-0.012}$ on 0/3 folds. Features the model happily splits on can still be actively
harmful once the comparison controls for what the rest of the feature set supplies.}

\finding{\textbf{Screening and full CV agree on rank but not magnitude.} \texttt{02\_price\_relative}
measured $+0.036$ screening versus $\approx+0.006$ full; physics $-0.012$ screening versus
$+0.001$ full. Rank ordering of the big movers is preserved, so screening is usable for
\emph{selection}, but screening deltas must never be quoted as results.}

\subsubsection*{Data-usage audit.}
With the score plateauing around 0.81--0.82, the question is what information remains untouched.

\paragraph{Used.} Price (raw, case-relative, forward/backward windows, shape); water value (raw,
system factor, per-reservoir deviation, terminal); storage (raw, seasonal percentile, and via
\texttt{vol\_head} to gross head and MWh/m\textsuperscript{3}); inflow (raw, trailing 30/90/365d,
case-horizon sums); topology (reservoir \texttt{max\_vol}/\texttt{lrl}/\texttt{vol\_head}, plant
\texttt{outlet\_line}, generator \texttt{p\_min}/\texttt{p\_max}/\texttt{startcost\_const}/
\texttt{penstock}, connections at 1-hop and full cascade).

\paragraph{Available and entirely unused.}
\begin{itemize}\itemsep1pt
  \item \texttt{Constraint\_min\_volume.csv} and \texttt{Constraint\_min\_flow.csv} --- two whole
        files, zero features built.
  \item 37 \texttt{river} records: \texttt{up\_head\_flow\_curve} (head $\to$ maximum flow, i.e.\
        hard throughput ceilings), \texttt{flow\_cost\_const} (spill penalties),
        \texttt{upstream\_elevation}.
  \item 5 \texttt{tunnel} records: \texttt{length}, \texttt{loss\_factor}, endpoint heights.
  \item Generator \texttt{turb\_eff\_curves} used only at their peak; the full flow-dependent
        curves, \texttt{gen\_eff\_curve} and \texttt{stopcost\_const} are unused.
  \item Plant \texttt{main\_loss} and \texttt{penstock\_loss} (head losses, quadratic in flow).
  \item The 144-hour overlap between consecutive cases, and case-to-case input deltas.
  \item Case-curve shape (PCA / Fourier of the 168h path).
\end{itemize}

\paragraph{Not in the dataset at all --- candidate explanations for the ceiling.}
The labels are binary ON/OFF, but SHOP solves for \emph{production levels}; we never see MW per
unit per hour, so two schedules with very different dispatch look identical to us. We also lack
the optimizer's internal state (dual values, marginal costs), any demand or reserve requirement,
and maintenance outages. One open question worth resolving: the inflow series is documented as
``historical'' --- if SHOP was driven by an inflow \emph{forecast} while we hold actuals, then our
inflow features are systematically the wrong quantity.

\subsubsection*{Next.} Prune to the smallest feature set that holds the score, then retrain that
set deeper. A smaller set that matches 299 features would confirm that most of round 2 was
redundancy rather than signal, and makes every subsequent experiment cheaper.
